\itemhead{2.3 思考 · 未编号题组（教材p.50）}{XJ-B1-C02-TH-P050-Y1033}
\label{XJ-B1-C02-TH-P050-Y1033}
{\footnotesize 湘教版教材p.50 / PDF第57页 / 2.3 思考 / 未编号题组}
\field{原题}{从 2.2 节可以知道，一元二次方程与相应的二次函数有着密切的联系，一元二次方程的根就是对应二次函数的零点。那么，能否利用它们之间的联系，借助二次函数求解一元二次不等式呢？}
\field{参考答案}{可以。先求二次函数零点，再根据开口和图象在 $x$ 轴上、下方的区间，分别得到函数值为正、负的解集；无零点和重根情形须分别讨论。}
{\small \field{关联知识}{I04 一元二次不等式与参数分类（教材帮p.64、65、66 / PDF73、74、75）；Q01 二次函数、零点、判别式与根系关系（教材帮p.65 / PDF74）}
\field{方法与易错点}{化标准式、查系数、求根或判别式、按图定号，最后与实际范围取交集。 易错：套“内小外大”却忘开口；丢重根或降次；漏整数产量限制。}
\field{答案依据}{答案册未提供本栏目的独立对应答案；依据原题、教材概念补写参考回应。开放性总结不设唯一数值答案。}}
